% ================================================================
%  conteudo/revisao.tex
% ================================================================

\section{Revisão da Literatura}

\subsection{Leveduras e fungos filamentosos}

As leveduras são microrganismos eucarióticos unicelulares pertencentes ao reino Fungi, caracterizados por sua capacidade de se reproduzir predominantemente por brotamento (gemulação) e, em algumas espécies, por fissão binária. Seu metabolismo é versátil, permitindo o crescimento tanto em condições aeróbias quanto anaeróbias, o que as torna organismos de grande relevância em processos biotecnológicos (VIEIRA; FERNANDES, 2012). Entre as espécies mais estudadas, \textit{Saccharomyces cerevisiae} destaca-se pela sua ampla aplicação industrial, sendo empregada na produção de etanol, pão, cervejas e vinhos (BERES, 2013).

Os fungos filamentosos, por sua vez, caracterizam-se por uma organização multicelular estruturada em hifas, que se ramificam formando o micélio. Essa complexidade estrutural permite-lhes colonizar diferentes substratos e secretar uma ampla variedade de enzimas extracelulares, como proteases, lipases e celulases, com elevado potencial de aplicação na indústria alimentícia, farmacêutica e de biocombustíveis (FARINAS; BARBOZA, 2012). Segundo Carvalho (2010), a capacidade de realizar modificações pós-traducionais complexas torna esses fungos plataformas vantajosas para a produção de proteínas recombinantes de interesse industrial.

\subsection{Microrganismos associados a frutas}

A superfície e a polpa das frutas constituem nichos ecológicos ricos em açúcares, vitaminas e minerais que favorecem o estabelecimento de comunidades microbianas diversificadas. As uvas, em particular, abrigam uma microbiota nativa composta majoritariamente por leveduras dos gêneros \textit{Hanseniaspora}, \textit{Candida}, \textit{Metschnikowia} e, em menor proporção, \textit{Saccharomyces}, cuja predominância aumenta no decorrer do processo fermentativo espontâneo (BERES, 2013). A composição dessa microbiota é fortemente influenciada pelo pH do fruto, pelo estágio de maturação e pelas condições climáticas da região de cultivo.

Conforme descrito pela EEEP (2012), frutas com pH inferior a 4,0 tendem a selecionar naturalmente leveduras e fungos filamentosos em detrimento de bactérias, visto que esses microrganismos apresentam maior tolerância à acidez. Essa seletividade é determinante para a definição das comunidades estabelecidas na superfície dos frutos e influencia diretamente o perfil de metabólitos produzidos durante fermentações espontâneas ou industriais.

\subsection{Meios de cultura sólidos}

O sucesso do isolamento laboratorial depende da escolha de suportes nutricionais que mimetizem o habitat natural do microrganismo. Os meios de cultura sólidos, estabilizados pelo uso do ágar, permitem a individualização de colônias e a análise macroscópica de suas características. Para o estudo de fungos e leveduras, utilizam-se formulações com alta densidade de carboidratos e pH ajustado, como o Ágar Batata Dextrose (BDA). De acordo com Soares, Ishimoto e Battestin (2013), o manejo desses meios é uma etapa estratégica na bioprospecção, servindo para enriquecer seletivamente linhagens com potencial para a secreção de enzimas e metabólitos de interesse comercial.

\subsection{Técnicas de isolamento microbiano}

O isolamento microbiano é o conjunto de procedimentos que visa a obtenção de culturas puras a partir de amostras mistas. A técnica de esgotamento por estrias em placa de Petri é um dos métodos mais difundidos, consistindo no arrastar do inóculo sobre a superfície do meio sólido de forma progressivamente diluída, de modo a promover a separação individual das células e o desenvolvimento de colônias isoladas (CÂMARA, 2023). Esse procedimento, aliado ao uso de materiais estéreis e à manutenção de um campo asséptico com auxílio de bico de Bunsen, é fundamental para evitar contaminações cruzadas.

Para amostras de superfície de frutos, a coleta por swab é amplamente utilizada por ser um método simples, não destrutivo e eficaz na recuperação da microbiota superficial. Após a inoculação, as placas são incubadas em estufas bacteriológicas a temperaturas controladas, geralmente entre 25 e 30~°C, por períodos que variam de 2 a 7 dias, dependendo do grupo microbiano de interesse (SOARES; ISHIMOTO; BATTESTIN, 2013).
